\documentclass[11pt]{article}

\usepackage{amsmath,amssymb,amsthm}
\usepackage[margin=1in]{geometry}

\newtheorem{proposition}{Proposition}
\newtheorem{lemma}{Lemma}
\newtheorem{remark}{Remark}

\newcommand{\C}{\mathbb C}
\newcommand{\R}{\mathbb R}
\newcommand{\1}{\mathbf 1}
\newcommand{\ddA}{\,dA}
\newcommand{\dz}{\partial_z}
\newcommand{\dbz}{\partial_{\bar z}}

\title{Section 4 in Lowest-Landau Language}
\author{}
\date{}

\begin{document}
\maketitle

\section{Purpose}

This note records the translation of the rigidity argument in Section 4 of
\texttt{main.tex} into the language of the magnetic Landau operator.  The main
point is that, for the usual Gaussian window used in the paper, the relevant
Landau level is the lowest Landau level, with magnetic field normalized by the
Fock weight
\[
        \|F\|_{\mathcal F^2}^2
        =\int_{\C}|F(z)|^2 e^{-\pi |z|^2}\ddA(z).
\]
Thus the constant which is \(2\) in the convention
\(e^{-|z|^2/2}\) becomes \(2\pi\) in the convention of the paper.

The conclusion is that Section 4 should be viewed as a magnetic Schwarz
potential problem at the bottom Landau energy.  The disk rigidity then follows
from the same mechanism as in the second proof of Section 4, provided one has
a magnetic analogue of the compact-support potential/divisibility step.

\section{Fock normalization and the Landau operator}

Set \(z=x+iy\), and choose the symmetric gauge
\[
        A_\pi(x,y)=(-\pi y,\pi x).
\]
Write
\[
        L_\pi:=(\nabla-iA_\pi)^2
        =(\partial_x+i\pi y)^2+(\partial_y-i\pi x)^2.
\]
This is the negative magnetic Laplacian.  The positive magnetic Laplacian is
\[
        H_\pi:=-L_\pi.
\]
Define the magnetic Cauchy--Riemann operators
\[
        Q_\pi:=2\dbz+\pi z,
        \qquad
        P_\pi:=2\dz-\pi\bar z .
\]
Then
\[
        Q_\pi=2e^{-\pi |z|^2/2}\dbz e^{\pi |z|^2/2},
        \qquad
        P_\pi=2e^{\pi |z|^2/2}\dz e^{-\pi |z|^2/2},
\]
and
\begin{equation}\label{eq:factorization}
        P_\pi Q_\pi=L_\pi+2\pi .
\end{equation}
Consequently,
\[
        Q_\pi \psi=0
        \quad\Longleftrightarrow\quad
        \psi(z)=F(z)e^{-\pi |z|^2/2},
        \qquad F \ \text{entire}.
\]
These are the lowest Landau states, and they satisfy
\[
        L_\pi\psi=-2\pi\psi,
        \qquad
        H_\pi\psi=2\pi\psi.
\]
If one rescales the field so that the Gaussian is \(e^{-|z|^2/2}\), then
\(2\pi\) is replaced by \(2\).  This is the source of the normalization shift.

\section{Localization as a lowest-Landau Toeplitz operator}

The Bargmann-Fock localization operator in the paper is
\[
        T_UF(w)=\int_U F(z)e^{\pi w\bar z}e^{-\pi |z|^2}\ddA(z).
\]
Under the unitary identification
\[
        F\longmapsto \psi_F:=F e^{-\pi |z|^2/2},
\]
this is the compressed multiplication operator
\[
        \psi_F\longmapsto \Pi_0(\1_U\psi_F),
\]
where \(\Pi_0\) is the orthogonal projection onto the lowest Landau level
\[
        \mathcal L_0:=\ker(H_\pi-2\pi)=\ker Q_\pi .
\]
Thus the Gaussian eigenfunction condition \(T_U1=\lambda\) is equivalent to
\[
        \Pi_0(\1_U\psi_0)=\lambda\psi_0,
        \qquad
        \psi_0=e^{-\pi |z|^2/2}.
\]
Testing against the reproducing kernels gives the weighted quadrature identity
\begin{equation}\label{eq:k0-quadrature}
        \int_U e^{-\pi |z|^2}e^{\pi w\bar z}\ddA(z)=\lambda,
        \qquad w\in\C .
\end{equation}
Equivalently,
\[
        \int_U e^{-\pi |z|^2}\bar z^m\ddA(z)=0,
        \qquad m\ge 1.
\]

More generally, if the Hermite function \(h_k\) is an eigenfunction, then its
Bargmann transform is a nonzero multiple of \(z^k\).  Section 4 obtains
\begin{equation}\label{eq:k-quadrature}
        \int_U |z|^{2k}e^{-\pi |z|^2}e^{\pi w\bar z}\ddA(z)
        =c_{k,\lambda},
        \qquad w\in\C ,
\end{equation}
after differentiating \(k\) times in \(w\).  This is still a lowest-Landau
quadrature identity.  The index \(k\) is angular momentum inside the lowest
Landau level; it is not a higher Landau level.

\section{The magnetic Schwarz-potential formulation}

The natural magnetic analogue of the logarithmic potential in Section 4 is a
compactly supported solution of
\begin{equation}\label{eq:magnetic-potential}
        (L_\pi+2\pi)u
        =
        |z|^{2k}e^{-\pi |z|^2/2}\1_U
        -
        c_{k,\lambda}\delta_0 .
\end{equation}
Indeed, if \(h(z)=G(z)e^{-\pi |z|^2/2}\) is in the lowest Landau level, then
pairing \eqref{eq:magnetic-potential} with the conjugate reproducing family
recovers exactly \eqref{eq:k-quadrature}.  Conversely, a compactly supported
solution of \eqref{eq:magnetic-potential} would be the Landau analogue of the
compactly supported logarithmic potential supplied in Section 4 by the
Fourier/divisibility argument.

If \(u\) is supported in \(\overline U\) and is sufficiently regular up to
\(\partial U\), extension by zero gives the overdetermined boundary conditions
\begin{equation}\label{eq:magnetic-boundary}
        u=0,
        \qquad
        \partial_\nu^A u:=\nu\cdot(\nabla-iA_\pi)u=0
        \quad\text{on }\partial U.
\end{equation}
Since the trace of \(u\) vanishes, its tangential magnetic derivative also
vanishes on \(\partial U\).  Hence the full magnetic gradient vanishes there,
and in particular
\[
        Q_\pi u=0
        \quad\text{on }\partial U .
\]

\section{Mirroring the holomorphic auxiliary function from Section 4}

Assume for the moment that \eqref{eq:magnetic-potential} has a compactly
supported solution satisfying \eqref{eq:magnetic-boundary}.  By
\eqref{eq:factorization},
\[
        P_\pi(Q_\pi u)
        =
        |z|^{2k}e^{-\pi |z|^2/2}\1_U
        -
        c_{k,\lambda}\delta_0 .
\]
Inside \(U\setminus\{0\}\), introduce the radial primitive
\[
        \Phi_k(t):=\frac12\int_0^t s^k e^{-\pi s}\,ds,
        \qquad t=|z|^2.
\]
Then
\[
        \Psi_k(z):=
        e^{\pi |z|^2/2}\frac{\Phi_k(|z|^2)}{\bar z}
\]
satisfies
\begin{equation}\label{eq:P-primitive}
        P_\pi\Psi_k
        =
        |z|^{2k}e^{-\pi |z|^2/2}
        \qquad (z\ne 0).
\end{equation}
Also,
\[
        P_\pi\left(e^{\pi |z|^2/2}\frac1{\bar z}\right)
        =
        2\pi\delta_0
\]
in the sense of distributions.  Therefore
\[
        G(z):=
        Q_\pi u(z)-\Psi_k(z)
        +\frac{c_{k,\lambda}}{2\pi}
        e^{\pi |z|^2/2}\frac1{\bar z}
\]
satisfies
\[
        P_\pi G=0
        \quad\text{in }U.
\]
Equivalently,
\[
        e^{-\pi |z|^2/2}G(z)
\]
is antiholomorphic in \(U\).  Multiplying by \(\bar z\), the singularity at the
origin is removable, and
\[
        W(z):=\bar z\,e^{-\pi |z|^2/2}G(z)
\]
is antiholomorphic in \(U\).

On \(\partial U\), the boundary condition gives \(Q_\pi u=0\), so
\[
        W|_{\partial U}
        =
        \frac{c_{k,\lambda}}{2\pi}
        -
        \Phi_k(|z|^2).
\]
This boundary value is real.  Hence \(\operatorname{Im}W\) is harmonic in \(U\),
continuous up to the boundary, and vanishes on \(\partial U\).  By the maximum
principle,
\[
        \operatorname{Im}W\equiv 0.
\]
Thus \(W\) is constant.  Restricting again to the boundary gives
\[
        \Phi_k(|z|^2)=\text{constant}
        \quad\text{on }\partial U.
\]
Since
\[
        \Phi_k'(t)=\frac12t^k e^{-\pi t}>0
        \qquad (t>0),
\]
the radius \(|z|\) is constant on \(\partial U\).  Under the same connectedness
and boundary hypotheses used in Section 4, this forces \(U\) to be a disk
centered at the origin.

\begin{proposition}[Conditional Landau-Schwarz rigidity]
Let \(U\subset\C\) be bounded, connected, and sufficiently regular, and assume
that \(0\in U\).  Suppose the quadrature identity
\eqref{eq:k-quadrature} gives rise to a compactly supported solution
\(u\) of \eqref{eq:magnetic-potential} with support in \(\overline U\) and
with boundary conditions \eqref{eq:magnetic-boundary}.  Then \(U\) is a disk
centered at the origin.
\end{proposition}

\section{Where the spectral nonresonance enters}

In the field normalization of this note, the bottom positive Landau energy is
\(2\pi\).  In the normalization of \texttt{landau.tex}, where the Gaussian is
\(e^{-|z|^2/2}\), it is \(2\).  The statement
\[
        2\pi\notin\sigma(H_{\pi,U}^D)
\]
for the Dirichlet magnetic Laplacian on a bounded domain is the corresponding
nonresonance statement.

Indeed, if \(H_{\pi,U}^D v=2\pi v\), then \(Q_\pi v=0\) by the factorization,
so
\[
        v(z)=F(z)e^{-\pi |z|^2/2}
\]
for some holomorphic \(F\).  The Dirichlet trace forces \(F=0\) on
\(\partial U\), hence \(F\equiv 0\) under the usual connected-domain
hypotheses.  Thus \(2\pi\) is not a Dirichlet eigenvalue.

This nonresonance is useful because it gives uniqueness and solvability for
the interior Dirichlet problem at the Landau energy.  It is not, by itself, the
disk rigidity.  The rigidity comes from the stronger free-boundary information:
the solution is compactly supported, so both \(u\) and the magnetic normal
derivative vanish on the boundary.

\section{What remains to make this a proof}

To turn the preceding conditional argument into a replacement proof of
Section 4, one needs the following inputs.

\begin{enumerate}
    \item \textbf{Magnetic divisibility / compact-support potential.}
    Starting from \eqref{eq:k-quadrature}, prove the existence of
    \(u\in\mathcal E'(\C)\), or preferably a function with enough Sobolev
    regularity, satisfying \eqref{eq:magnetic-potential}.  This is the magnetic
    analogue of the Fourier/Paley--Wiener divisibility step used in Section 4
    to obtain \(\Delta u=f\1_U-c\delta_0\).

    \item \textbf{Support identification.}
    Show that the compactly supported solution may be chosen so that
    \(\operatorname{supp}u\subset\overline U\), or prove directly that it
    vanishes in the exterior component.  This is what turns the PDE into a
    free-boundary problem.

    \item \textbf{Boundary regularity and transmission.}
    Establish enough regularity up to \(\partial U\) to justify
    \(u=0\) and \(\partial_\nu^A u=0\) on the boundary after extension by zero.
    Once these are available, the first-order argument above is purely local.

    \item \textbf{Treatment of the origin.}
    As in Section 4, one must show that \(0\in U\) and \(0\notin\partial U\),
    or otherwise remove a small centered disk and adjust the constant.  The
    auxiliary function uses \(1/\bar z\), so this point has to be handled
    explicitly.

    \item \textbf{Topology.}
    The boundary conclusion is that \(\partial U\) lies on a centered circle.
    The final statement that \(U\) is the disk uses the same connectedness or
    simply-connectedness assumptions as Section 4.
\end{enumerate}

\section{Summary}

The Gaussian-window Section 4 problem is a lowest-Landau-level inverse problem:
\[
        T_U z^k=\lambda z^k
        \quad\Longleftrightarrow\quad
        \Pi_0(\1_U z^k e^{-\pi |z|^2/2})
        =
        \lambda z^k e^{-\pi |z|^2/2}.
\]
After differentiating in the reproducing parameter, this becomes the weighted
quadrature identity \eqref{eq:k-quadrature}.  The corresponding magnetic
Schwarz potential is \eqref{eq:magnetic-potential}, at the Landau energy
\(2\pi\) for the positive operator \(H_\pi\).

If the missing compact-support potential step is supplied, then the rest of
Section 4 has a direct Landau analogue: factor \(L_\pi+2\pi=P_\pi Q_\pi\),
subtract the explicit radial first-order primitive \(\Psi_k\), obtain an
antiholomorphic auxiliary function with real boundary values, and conclude that
\(|z|\) is constant on \(\partial U\).  This recovers the disk.

\end{document}
